\title{Direct Preference Optimization: Your Language Model is Secretly a Reward Model}

\begin{abstract}
Although large unsupervised language models (LMs) acquire extensive knowledge about the world and certain reasoning abilities, exerting fine-grained control over their outputs remains challenging due to the inherently unsupervised nature of their training process. Current approaches aiming to enhance model steerability typically rely on gathering human judgments regarding the relative quality of generated outputs, subsequently fine-tuning the unsupervised LM to better reflect these preferences, most often through reinforcement learning from human feedback (RLHF). Nevertheless, RLHF involves a complicated and frequently unstable pipeline: initially fitting a reward model to human preference data, and then optimizing the large language model via reinforcement learning to increase this proxy reward, while constraining updates to maintain proximity to the initial model. In this work, we propose a novel reward model parameterization within the RLHF framework, which permits closed-form computation of the associated optimal policy. This innovation enables us to reframe the standard RLHF objective as a straightforward classification loss. The proposed method, termed \textit{Direct Preference Optimization} (DPO), is efficient, robust, and requires minimal computational resources, removing the necessity for LM sampling during fine-tuning and extensive hyperparameter adjustment. Experimental results demonstrate that DPO fine-tunes LMs to reflect human preferences on par with or superior to conventional techniques. Specifically, DPO surpasses PPO-based RLHF in controlling sentiment in generated outputs, and attains comparable or improved performance in tasks like summarization and single-turn dialogue—while being significantly easier to implement and train.
\end{abstract}

\section{Introduction}
Large unsupervised language models (LMs), when exposed to extensive datasets, demonstrate remarkable and unexpected abilities~\citep{chowdhery2022palm, brown2020language, touvron2023llama,bubeck2023sparks}. Nonetheless, these models are trained on data produced by humans whose objectives, levels of expertise, and motivations differ widely. Not all of these human traits are desirable for imitation; for instance, although it is beneficial for an AI coding assistant to \textit{recognize} prevalent programming errors in order to fix them, we nevertheless prefer the model’s code generation to reflect the (sometimes infrequent) superior programming practices present in its training sources. Analogously, we may prefer that a language model is \textit{informed} about a widespread misconception—such as one endorsed by half the population—but we do not wish the model to propagate this falsehood in half of its answers concerning the topic. Thus, it is imperative to differentiate and prioritize the model’s \emph{intended outputs and conduct} from the broad expanse of its \textit{knowledge and skills}, a foundational step for creating AI that is reliable, high-performing, and manageable \citep{ouyang2022training}. Existing approaches generally align LMs to human values and preferences via reinforcement learning (RL), but as we will demonstrate, the objective optimized by these RL-based methods can be directly realized through a straightforward binary cross-entropy loss, streamlining the entire preference alignment process.

Broadly speaking, current techniques instill preferred behaviors in language models by leveraging collections of human-annotated preferences, which serve as a template for behaviors deemed beneficial or safe by people. This phase of preference alignment typically follows massive unsupervised pre-training on extensive text corpora. Although one intuitive method for achieving such alignment is to perform supervised fine-tuning with high-quality human example responses, the most prominent set of techniques involve reinforcement learning from human (or AI) feedback (RLHF/RLAIF; \citep{christiano2017deep,bai2022constitutional}). In RLHF, a reward model is trained using preference data provided by humans, and reinforcement learning is then employed to adjust the model policy so as to favor responses with higher reward scores, while simultaneously restraining deviation from the initial model parameters. Despite the fact that RLHF approaches yield models with superior conversational and programming skills, the process is substantially more intricate than supervised fine-tuning; it requires the training of several language models and necessitates on-policy sampling during training, thus imposing considerable computational burdens.

This work introduces a method for optimizing language models to better align with human preferences without relying on explicit reward models or reinforcement learning frameworks. We present \textit{{Direct Preference Optimization} (DPO)}, an algorithm that implicitly targets the same objective as common RLHF methods—specifically, reward maximization under a KL-divergence regularization—but offers a more accessible implementation and training process. Conceptually, DPO modifies the model by boosting the log odds of preferred responses over less favored alternatives, while applying a per-sample dynamic importance weight. This weighting is crucial for mitigating model collapse, which we observe if one simply maximizes the probability ratio between preferred and dispreferred outputs. DPO builds on a theoretical basis provided by preference models (such as the Bradley-Terry model; \cite{bradley1952rankanalysis}) to assess the alignment between proposed reward functions and real-world preference annotations. Distinct from other approaches that leverage these models to supervise a reward function—which is subsequently used to train a policy—DPO employs a variable transformation that directly expresses the preference loss in terms of the policy itself. Consequently, with a dataset of human-annotated preferences regarding model outputs, DPO enables direct policy optimization through a straightforward binary cross-entropy loss, yielding a policy that is optimal relative to an implicit reward function derived from the preference data.

Our principal contribution is the development of Direct Preference Optimization (DPO), a straightforward algorithm for preference-based language model training that dispenses with reinforcement learning. Experimental results indicate that DPO matches or exceeds the performance of existing preference-based approaches, including RLHF methods using PPO, across various tasks such as sentiment control, summarization, and conversational modeling, on models containing up to 6 billion parameters.

As self-supervised language models have grown in scale, they have demonstrated an ability to perform certain tasks without explicit task-specific training, either by responding to zero-shot prompts \citep{radford2019language} or by utilizing few-shot examples \citep{gpt3,megatron,chowdhery2022palm}. Nonetheless, these models exhibit much stronger performance on downstream applications and improved alignment with user intentions when fine-tuned with datasets containing explicit instructions paired with human-generated responses \citep{mishra-etal-2022-cross,sanh2022multitask,chung2022scaling,thoppilan2022lamda}. This process, referred to as 'instruction-tuning', equips LLMs to better extrapolate to novel instructions beyond the original training set, thereby enhancing their practical utility \citep{chung2022scaling}. Even with these gains, collecting \textit{relative} human ratings of response quality is typically less demanding than gathering high-quality expert-written samples, leading to subsequent methods that leverage datasets of human preferences to further fine-tune LLMs. This approach has yielded advances in domains such as translation \citep{kreutzer-etal-2018-reliability}, summarization \citep{stiennon2022learning,ziegler2020finetuning}, narrative generation \citep{ziegler2020finetuning}, and complex instruction-following tasks \citep{ouyang2022training,ramamurthy2023is}. Typically, these systems start by training a neural reward model to align with the collected preference data under a ranking model like the Bradley-Terry model \citep{bradley1952rankanalysis}, and subsequently update the language model parameters to maximize the predicted reward via reinforcement learning strategies such as REINFORCE \citep{williams1992reinforce}, proximal policy optimization (PPO; \cite{schulman2017proximal}), or related variants \citep{ramamurthy2023is}. Another closely-aligned research direction employs instruction-following LLMs, refined using human feedback, to autonomously generate synthetic preference data targeting attributes like safety or harmlessness \citep{bai2022constitutional}, with human involvement limited to providing textual rubrics guiding the model’s annotations. These approaches bring together two major research areas: reinforcement learning for language model training on diverse objectives~\citep{Ranzato2015SequenceLT,paulus2018a,wu2018learning}, and frameworks for learning from human evaluative feedback \citep{christiano2017deep,kupcsik2018learning}. Nevertheless, although training with relative preferences is appealing, applying reinforcement learning to fine-tune large language models remains operationally challenging. In this context, our work introduces a theoretically-grounded method for optimizing with relative preferences that circumvents the need for reinforcement learning.

Research on learning policies based on preferences, independent of language applications, has been explored within both bandit and reinforcement learning frameworks, leading to the development of multiple methodologies. In contextual bandit scenarios, preference or ranking feedback—rather than explicit reward signals—is referred to as the contextual dueling bandit (CDB; \cite{yue2012karmed,dudik2015contextual}). When explicit rewards are not available, theoretical work on CDBs replaces the traditional definition of optimality with the concept of a \textit{von Neumann winner}: a policy that achieves at least a 50\% expected win rate against any other candidate policy \citep{dudik2015contextual}. Importantly, in CDBs, feedback is provided in an online fashion, whereas, in the context of learning from human preferences, algorithms are often trained on a static dataset composed of pre-collected, preference-annotated action pairs \citep{yan2022human}. In a similar vein, \textit{preference-based RL} (PbRL) learns from binary comparisons drawn from an \textit{unknown} underlying scoring function in lieu of reward feedback \citep{BusaFekete2014,ruiz2023dueling}. Numerous PbRL methods exist—including those designed to leverage previously collected (off-policy) preference data—yet these generally proceed by first estimating a latent scoring (reward) function, and then optimizing the policy with respect to this estimate \citep{jain2013learning,BusaFekete2014,christiano2017deep,sadigh2017active,kupcsik2018learning}. In contrast, we introduce a unified policy learning strategy that directly adjusts the policy to conform to the observed preferences.

\section{Preliminaries}\label{section:prelims}

We briefly outline the RLHF pipeline as described by \citeauthor{ziegler2020finetuning}, with further elaborations in subsequent works \citep{stiennon2022learning, bai2022training, ouyang2022training}. The process is typically divided into three stages: (1) supervised fine-tuning (SFT), (2) collecting preference data and training a reward model, and (3) reinforcement learning-based optimization.

\textbf{SFT}: The RLHF procedure typically starts by adapting a pre-trained language model (LM) through supervised learning, using curated high-quality datasets targeted at specific downstream tasks (such as dialogue or summarization), resulting in a fine-tuned model denoted by $\pi^\text{SFT}$.

\textbf{Reward Modelling Phase}: In the following stage, the fine-tuned SFT model is supplied with prompts $x$ to generate response pairs $(y_1, y_2)\sim \pi^\text{SFT}(y \mid x)$. These output pairs are presented to human annotators, who indicate their preferred response—formally, $y_w\succ y_l \mid x$, where $y_w$ and $y_l$ represent the favored and unfavored answers among $(y_1, y_2)$, respectively. The underlying assumption is that these preferences arise from some hidden reward function $r^*(y, x)$, to which we do not have direct access. Preference modeling can be implemented via various techniques, with the Bradley-Terry (BT) \cite{bradley1952rankanalysis} model being especially prevalent (though more comprehensive Plackett-Luce models \citep{plackett1975analysis, luce2012individual} are also suitable if multiple ranked completions are available). According to the BT model, the distribution of human preferences $p^*$ is formalized as:

\begin{equation}\label{eq:bradley-terry}
    p^*(y_1\succ y_2 \mid x)=\frac{\exp\left(r^*(x, y_1)\right)}{\exp\left(r^*(x, y_1)\right) + \exp\left(r^*(x, y_2)\right)}.
\end{equation}

Given a fixed set of comparative samples $\mathcal{D}=\bigl\{x^{(i)}, y_w^{(i)}, y_l^{(i)}\bigr\}_{i=1}^N$ drawn according to $p^*$, we can introduce a parameterized reward function $r_{\phi}(x, y)$ and fit the model parameters using the maximum likelihood principle. Casting the learning objective as a binary classification task, the associated negative log-likelihood loss is expressed as:

\begin{equation}\label{eq:reward_model}
    \mathcal{L}_R(r_{\phi}, \mathcal{D}) = -\mathbb{E}_{(x, y_w, y_l)\sim \mathcal{D}}\bigl[\log \sigma(r_{\phi}(x, y_w)- r_{\phi}(x, y_l))\bigr]
\end{equation}

where $\sigma$ denotes the logistic sigmoid. For language models, the function $r_{\phi}(x, y)$ is typically initialized from the SFT model $\pi^\text{SFT}(y \mid x)$ and augmented with a linear output layer atop the final transformer block, which produces a scalar reward estimate \cite{ziegler2020finetuning}. To stabilize training and reduce reward variance, it is common practice to normalize the outputs of $r_\phi$, enforcing $\mathbb{E}_{x,y\sim \mathcal{D}}\left[r_\phi(x, y)\right] = 0$ for each $x$.

\textbf{RL Fine-Tuning Phase}: In the reinforcement learning stage, the previously trained reward model acts as a feedback signal for the language model. Following existing literature~\citep{jaques2017sequence, jaques2020human}, the training objective takes the form:

\begin{equation}\label{eq:RL}
\max_{\pi_{\theta}}  \mathbb{E}_{x\sim \mathcal{D}, y\sim \pi_{\theta}(y \mid x)}\bigl[r_{\phi}(x, y)\bigr] - \beta\mathbb{D}_{\textrm{KL}}\bigl[\pi_{\theta}(y\mid x)\mid \mid \pi_\text{ref}(y\mid x)\bigr],
\end{equation}

where the hyperparameter $\beta$ regulates the regularization strength against a reference policy $\pi_\text{ref}$, typically set to the pre-trained SFT model $\pi^\text{SFT}$. In application, the policy $\pi_\theta$ is also initialized from $\pi^\text{SFT}$. This constraint is necessary to ensure that the policy remains within a region where the reward model is reliable, while also encouraging diverse generation and avoiding collapse to narrow high-reward solutions. Because text generation is a discrete process, this objective cannot be directly differentiated and is therefore commonly approached with reinforcement learning algorithms. The prevailing method \citep{ziegler2020finetuning, stiennon2022learning, bai2022training, ouyang2022training} involves defining the modified reward $r(x, y) = r_{\phi}(x, y) -\beta (\log \pi_{\theta}(y\mid x) - \log \pi_\text{ref}(y\mid x))$ and optimizing it through PPO \cite{schulman2017proximal}. 

\section{Direct Preference Optimization}\label{sec:DPO}

Driven by the limitations of using reinforcement learning at scale for tuning language models, we aim to develop a more straightforward strategy for optimizing policies using preference data. In contrast to traditional RLHF pipelines, which rely on first fitting a reward function and then maximizing it via RL, our method builds upon a reward model parameterization that allows us to directly recover the optimal policy in closed form, thereby circumventing the need for RL-based optimization. As elaborated below, the central idea is to use a closed-form relationship between reward functions and their corresponding optimal policies, allowing us to reformulate the learning objective from one over reward functions to one defined on policy parameters. This variable substitution approach eliminates the necessity to explicitly fit a separate reward network, yet still operates under standard models of human preferences like the Bradley-Terry framework. In summary, the policy network is responsible for modeling both the language generation

\textbf{Reformulation of the DPO objective.} We begin with the conventional reinforcement learning objective described in Eq.~\ref{eq:RL}, defined for an arbitrary reward function $r$. Building on established literature~\citep{peters2007reinforcement, peng2019advantage, korbak2022reinforcement, go2023aligning}, it can be readily demonstrated that the optimal policy for the KL-regularized reward maximization problem in Eq.~\ref{eq:RL} assumes the following structure:

\begin{equation}\label{eq:op_policy}
    \pi_r(y\mid x) = \frac{1}{Z(x)}\pi_\text{ref}(y\mid x)\exp\left(\frac{1}{\beta}r(x, y)\right),
\end{equation}

where $Z(x) =\sum_{y}\pi_\text{ref}(y\mid x)\exp\left(\frac{1}{\beta}r(x, y)\right)$ serves as the partition function. The complete derivation is provided in Appendix \ref{app:derivation1}. Notably, even when the reward function $r^*$ is approximated via its maximum likelihood estimator $r_{\phi}$, computation of $Z(x)$ remains computationally prohibitive~\citep{korbak2022reinforcement, go2023aligning}, limiting the direct utility of this form. Nevertheless, it is possible to manipulate Eq.~\ref{eq:op_policy} to solve for the reward function in terms of the optimal policy $\pi_r$, reference policy $\pi_\text{ref}$, and the (generally intractable) partition function $Z(\cdot)$. To do so, we take the logarithm of both sides and rearrange, yielding:

\begin{equation}\label{eq:main_eq}
    r(x,y) =\beta \log \frac{\pi_r(y\mid x)}{\pi_\text{ref}(y\mid x)} + \beta \log Z(x).
\end{equation}

This alternative parameterization can be used for the ground-truth reward $r^*$ together with its induced optimal policy $\pi^*$. Importantly, the Bradley-Terry model is only sensitive to reward differences between pairs of outputs, i.e., ${p^*(y_1 \succ y_2 \mid x) = \sigma(r^*(x, y_1) - r^*(x, y_2))}$. Substituting the result from Eq.~\ref{eq:main_eq} for $r^*(x,y)$ into the Bradley-Terry preference formula Eq.~\ref{eq:bradley-terry}, the partition function $Z(x)$ cancels, permitting us to express the preference probability solely in terms of $\pi^*$ and $\pi_\text{ref}$. Thus, under the Bradley-Terry framework, the optimal RLHF policy $\pi^*$ conforms to:

\begin{equation}\label{eq:objective}
    p^*(y_1\succ y_2 \mid x)=\frac{1}{1 + \exp\left(\beta \log \frac{\pi^*(y_2\mid x)}{\pi_\text{ref}(y_2\mid x)} - \beta \log \frac{\pi^*(y_1\mid x)}{\pi_\text{ref}(y_1\mid x)}\right)}
\end{equation}

Appendix~\ref{app:derivation2} contains the detailed derivation. While this result leverages the Bradley-Terry assumption, analogous forms can be derived for the broader family of Plackett-Luce models~\citep{plackett1975analysis, luce2012individual}, as elaborated in Appendix~\ref{app:plackett_luce_models}.

Having related the likelihood of human preference data to the optimal policy rather than the underlying reward, we are able to specify a maximum likelihood estimation problem for a parameterized policy $\pi_\theta$. Mirroring the approach for learning reward models (cf. Eq.~\ref{eq:reward_model}), we arrive at the following policy objective:

\begin{equation}\label{eq:optimum_model}
    \mathcal{L}_\text{DPO}(\pi_{\theta}; \pi_\text{ref}) = -\mathbb{E}_{(x, y_w, y_l)\sim \mathcal{D}}\left[\log \sigma \left(\beta \log \frac{\pi_{\theta}(y_w\mid x)}{\pi_\text{ref}(y_w\mid x)} - \beta \log \frac{\pi_{\theta}(y_l\mid x)}{\pi_\text{ref}(y_l\mid x)}\right)\right].
\end{equation}

In this approach, we estimate an implicit reward via a different parameterization, such that the corresponding optimal policy is exactly $\pi_\theta$. Additionally, our method can be viewed as fitting a reparameterized Bradley-Terry model, thereby inheriting established theoretical guarantees such as consistency under appropriate assumptions on the distribution of preference data \cite{bong2022generalized}. Further theoretical analysis of DPO and its connections to related methods can be found in Section~\ref{sec:theory}.

\textbf{How does the DPO update operate?} To gain a mechanistic perspective on DPO, we examine the gradient of the loss function $\mathcal{L}_\text{DPO}$. Specifically, the derivative with respect to $\theta$ is given by:

\begin{multline*}\label{eq:gradient}
    \nabla_\theta \mathcal{L}_\text{DPO}(\pi_\theta;\pi_\text{ref}) = \\ -\beta\mathbb{E}_{(x, y_w, y_l) \sim \mathcal{D}} \bigg[\underbrace{\sigma(\hat{r}_\theta(x, y_l) - \hat{r}_\theta (x, y_w))}_\text{assigns larger weight when reward estimate is inaccurate}\bigg[\underbrace{\nabla_\theta\log \pi(y_w \mid x)}_\text{increase probability of $y_w$} - \underbrace{\nabla_\theta\log\pi(y_l \mid x)}_\text{decrease probability of $y_l$}\bigg]\bigg],
\end{multline*}

with $\hat{r}_\theta(x, y) = \beta \log \frac{\pi_\theta(y \mid x)}{\pi_\text{ref}(y \mid x)}$ serving as the reward implicitly specified by the parameterized language model $\pi_\theta$ in relation to the reference model $\pi_\text{ref}$ (elaborated upon in Section~\ref{sec:theory}). Conceptually, the gradient of $\mathcal{L}_\text{DPO}$ increases the probability of generating the more desirable (preferred) completion $y_w$, while suppressing the likelihood of the less desirable (dispreferred) completion $y_l$. Notably, the influence of each data point is modulated according to the extent that the implicit reward model $\hat{r}_\theta$ overvalues the dispreferred completion compared to the preferred one, scaled by $\beta$, reflecting both the error in ranking and the intensity of the KL constraint. Our empirical results indicate that this example-dependent reweighting is crucial—removing it (i.e., using a simpler version without the weighting term) can cause model degeneration, as evidenced in Appendix Table~\ref{tab:unlikelihood_generations}.

\textbf{DPO procedure.} The typical DPO workflow consists of the following steps: 1) For each prompt $x$, generate response samples $y_1, y_2 \sim \pi_\text{ref}(\cdot \mid x)$ and annotate them with human preference data to form the preference dataset $\mathcal{D} = \{x^{(i)}, y_w^{(i)}, y_l^{(i)}\}_{i=1}^N$; 2) Train the language model $\pi_\theta$ by minimizing the loss $\mathcal{L}_\text{DPO}$ using the specified $\pi_\text{ref}$, $\mathcal{D}$, and the chosen value of $\beta$. 
In applied settings, it is preferable to utilize existing public preference datasets instead of creating new samples and collecting human feedback from scratch. Since these datasets are generally obtained with $\pi^\text{SFT}$, we set $\pi_\text{ref} = \pi^\text{SFT}$ when possible. If $\pi^\text{SFT}$ is not accessible, we instead initialize $\pi_\text{ref}$ by maximizing the likelihood on the set of preferred completions ${(x, y_w)}$; specifically, ${\pi_\text{ref} = \argmax_{\pi}\mathbb{E}_{x, y_w \sim \mathcal{D}}\left[\log \pi(y_w \mid x)\right]}$. This approach addresses potential distribution mismatch between the unavailable true reference distribution and the $\pi_\text{ref}$ utilized in DPO. Additional information about implementation and chosen hyperparameters appears in Appendix~\ref{app:implementation}.

\section{Theoretical Analysis of DPO} Here, we further interpret the DPO framework, provide supporting theoretical results, and connect the strengths of DPO to the known limitations of actor-critic RLHF algorithms like PPO~\cite{schulman2017proximal}.

\label{sec:theory}

\subsection{Your Language Model Is Secretly a Reward Model} DPO eliminates the need for explicit reward modeling and separate RL-based policy training by employing a unified maximum likelihood objective. Notably, the objective in Eq. \ref{eq:main_eq} corresponds to the Bradley-Terry model, parameterized by the reward $r^*(x, y) = \beta \log\frac{\pi^*_\theta(y \mid x)}{\pi_\text{ref}(y \mid x)}$, and training our parametric model $\pi_{\theta}$ is thus equivalent to reward model fitting as in Eq. \ref{eq:reward_model}, under an appropriate variable transformation. In the following, we develop the theoretical basis for this reparameterization, demonstrate that it places no additional restriction on the set of learnable reward models, and show that it admits exact recovery of the optimal policy. We begin by introducing an equivalence relation over reward functions.

\begin{definition}
Two reward functions $r(x, y)$ and $r'(x, y)$ are defined to be equivalent if and only if there exists some function $f$ such that ${r(x, y)-r'(x, y) = f(x)}$.
\end{definition}

This equivalence relation is straightforward to verify, as it naturally partitions the collection of reward functions into equivalence classes. This leads us to the following lemmas:

\begin{lemma}\label{lemma:same_prefrence}
Within the Plackett-Luce framework—including the Bradley-Terry model—any two reward functions belonging to the same equivalence class produce identical preference distributions.
\end{lemma}

\begin{lemma}\label{lemma:same_policy}
Reward functions that fall within the same equivalence class will yield the same optimal policy when considering the constrained RL setting.
\end{lemma}

The proofs are elementary and can be found in Appendix \ref{app:lemma1}. The first lemma highlights a classic under-identification phenomenon associated with the Plackett-Luce family \cite{plackett1975analysis}. Because of this, it is common practice to enforce additional identifiability constraints in order to ensure meaningful MLE solutions for Eq. \ref{eq:reward_model} \cite{bong2022generalized}. The second lemma indicates that the optimal policy does not depend on the specific choice of reward function within an equivalence class, so for our overall objective, it suffices to recover any representative from the appropriate class. We establish the following theorem, with proof detailed in Appendix~\ref{app:thm1}:

\begin{theorem}\label{thm:main}
    Under certain mild conditions, any equivalence class of reward functions compatible with the Plackett-Luce (and, as a special case, the Bradley-Terry) models admits a parameterization of the form ${r(x, y) = \beta \log \frac{\pi(y\mid x)}{\pi_\text{ref}(y\mid x)}}$ for a suitable model $\pi(y\mid x)$ and a specified reference model $\pi_\text{ref}(y \mid x)$.
\end{theorem}

\begin{sproof}
    Let $r(x, y)$ be any reward function, and denote the associated optimal model as $\pi_r(y \mid x)$ (see Eq. \ref{eq:op_policy}). We show that a reward function in the equivalence class of $r$ can always be expressed through the proposed reparameterization. Define the transformation $f$ as
\begin{equation}
    f(r; \pi_\text{ref}, \beta)(x, y) = r(x, y) - \beta\log\sum_{y}\pi_\text{ref}(y\mid x)\exp\left(\frac{1}{\beta}r(x, y)\right)
\end{equation}
This operator $f$ normalizes $r$ via the logarithm of the partition function induced by $\pi_r$. Since the normalization depends solely on $x$, $f(r; \pi_\text{ref}, \beta)(x, y)$ belongs to the equivalence class determined by $r(x, y)$. Substituting $r$ with the right-hand side of Eq.~\ref{eq:main_eq}—valid for any reward—we have $f(r; \pi_\text{ref}, \beta)(x, y) = \beta \log \frac{\pi_r(y\mid x)}{\pi_\text{ref}(y\mid x)}$. Thus, $f$ yields a canonical form within the same equivalence class, confirming that the proposed reparameterization is fully expressive for our reward modeling purposes.
\end{sproof}

An alternative interpretation of Theorem~\ref{thm:main} is that it uniquely identifies which representative of each reward function equivalence class is chosen under the DPO reparameterization; specifically, it selects the reward function that satisfies

\begin{equation}\label{eq:lag_p}
     \sum_{y}\underbrace{\pi_\text{ref}(y\mid x)\exp\left(\frac{1}{\beta}r(x, y)\right)}_{=\pi(y\mid x)\text{, using Thm.~\ref{thm:main} reparam.}} = 1,
\end{equation}

which ensures that $\pi(y\mid x)$ forms a proper probability distribution (i.e., it is non-negative and sums to one). By comparing this with Eq.~\ref{eq:op_policy}, it becomes evident that Eq.~\ref{eq:lag_p} defines the partition function corresponding to the optimal policy arising from $r(x, y)$. The central idea in the DPO algorithm is to apply specific constraints to the otherwise ambiguous Plackett-Luce (and, notably, the Bradley-Terry) family of preference models. This approach retains the set of reward models that can be represented while, crucially, allowing the optimal policy in Eq.~\ref{eq:op_policy} to admit a closed-form solution for any prompt $x$.

\subsection{Instability of Actor-Critic Algorithms} Our framework also provides a means to analyze and understand sources of instability in commonly used actor-critic methods for RLHF, such as PPO. Adhering to the RLHF pipeline, we consider the RL fine-tuning procedure as described in Section~\ref{section:prelims}, and link it to the control-as-inference paradigm \cite{levine2018reinforcement} associated with the constrained RL formulation in \ref{eq:RL}. We consider a parameterized policy $\pi_{\theta}(y\mid x)$, and aim to minimize the divergence $\mathbb{D}_{\text{KL}}[\pi_{\theta}(y|x) \mid \mid \pi^*(y\mid x)]$ where $\pi^*$ denotes the optimal policy from Eq.~\ref{eq:optimum_model}, determined by the reward function $r_{\phi}(y, x)$. Through a series of algebraic steps, this results in the following objective:

\begin{equation}\label{eq:AC}
    \max_{\pi_{\theta}}\mathbb{E}_{\pi_{\theta}(y\mid x)}\bigg[\underbrace{r_{\phi}(x, y) -\beta\log\sum_{y}\pi_\text{ref}(y\mid x)\exp\left(\frac{1}{\beta}r_{\phi}(x, y)\right)}_{f(r_{\phi}, \pi_\text{ref}, \beta)} - \underbrace{\beta\log\frac{\pi_{\theta}(y\mid x)}{\pi_\text{ref}(y\mid x)}}_{\text{KL}}\bigg]
\end{equation}

This objective aligns with those previously optimized in the literature \citep{ziegler2020finetuning, stiennon2022learning, bai2022training, ouyang2022training}, where the DPO-equivalent reward is employed for the reward function class $r_{\phi}$. Within this framework, the normalization component in $f(r_{\phi}, \pi_\text{ref}, \beta)$ can be interpreted as the soft value function corresponding to the reference policy $\pi_\text{ref}$. Although the presence of this term does not influence the optimal policy, omitting it may lead to high variance in policy gradient estimates, thus destabilizing training. While it is possible to compensate for this normalization using an auxiliary learned value function, optimizing such a function presents its own challenges. An alternative approach adopted in earlier works involves normalizing rewards by referencing a human completion baseline—effectively approximating the normalization term with a single-sample Monte Carlo estimate. In contrast, the DPO reparameterization naturally results in a reward function that obviates the need for explicit baseline corrections.

\section{Experiments}
Here, we provide empirical assessments of DPO's effectiveness for directly training policies from preference data. We begin with a controlled text-generation scenario, investigating how efficiently DPO balances reward maximization against minimizing KL divergence from the reference policy, in comparison with prevalent preference-learning algorithms like PPO. Subsequently, we assess DPO on more demanding settings, involving larger models and challenging RLHF tasks such as summarization and dialogue. Our experiments show that, with minimal hyperparameter tuning, DPO generally matches or surpasses the performance of competitive baselines including RLHF with PPO, as well as the best-of-$N$ trajectory selection under a learned reward model. Prior to discussing these findings, we outline the experimental methodology; further implementation specifics can be found in Appendix~\ref{app:exp_details}.

\textbf{Tasks.} Our study investigates three distinct open-ended text generation tasks. Across all tasks, methods learn a policy based on a dataset of preferences denoted as $\mathcal{D}=\bigl\{x^{(i)}, y_w^{(i)}, y_l^{(i)}\bigr\}_{i=1}^N$. In the case of \textbf{controlled sentiment generation}, $x$ corresponds to the start of a movie review sourced from the IMDb dataset \cite{maas-EtAl:2011:ACL-HLT2011}, with the objective that the policy generates $y$ exhibiting positive sentiment. For a systematic assessment, preference pairs for this task are synthetically generated by applying a pre-trained sentiment classifier, ensuring that $p(\text{positive}\mid x,y_w)>p(\text{positive}\mid x,y_l)$. For SFT, GPT-2-large is fine-tuned to convergence using training reviews from the IMDb dataset (for more details, see App~\ref{app:sentiment_details}). The \textbf{summarization} task uses $x$ as a Reddit forum post, requiring the policy to generate $y$ summarizing the main ideas of the post. Building on established work, we utilize the Reddit TL;DR summarization dataset \citep{volske-etal-2017-tl} and human preference annotations from \citeauthor{stiennon2022learning}. Here, the SFT model is trained on summaries authored by humans for forum posts,\footnote{\url{https://huggingface.co/CarperAI/openai_summarize_tldr_sft}} and RLHF is performed via the TRLX \citep{leandro_von_werra_2023_7790115} framework. The human preference annotations were collected by \citeauthor{stiennon2022learning} from generations of a related SFT model. In the final task, \textbf{single-turn dialogue}, $x$ is a prompt from a human user, which may vary from factual queries to personal advice-seeking. Here, the policy must provide a relevant and engaging response $y$. We employ the Anthropic Helpful and Harmless dialogue dataset \citep{bai2022training}, consisting of 170,000 human-assistant dialogues, each culminating in two model-generated responses alongside a label for the preferred reply. As there is no publicly available pre-trained SFT model for this task, we construct one by fine-tuning a standard language model exclusively on the preferred responses.

\textbf{Evaluation.} We employ two distinct evaluation strategies in our experiments. To assess how well each algorithm performs in the context of constrained reward maximization, we analyze the trade-off between reward and KL-divergence from the reference policy within the controlled sentiment generation scenario; here, the full reward-KL frontier can be determined due to our access to the oracle reward signal (the sentiment classifier). However, since a ground-truth reward function is unavailable in real-world applications, we instead measure algorithmic performance by their \textit{win rate} when pitted against a baseline policy, using GPT-4 as a stand-in for human judges to evaluate summary quality and single-turn dialogue helpfulness in the summarization and dialogue tasks, respectively. For the summarization task, the gold-standard reference summaries in the test set serve as the baseline, while for dialogue, the test set’s preferred responses are used. Although prior work indicates that large language models can serve as more reliable automatic evaluators than conventional metrics \citep{Chen2023ExploringTU}, we further conduct a human evaluation, presented in Sec.~\ref{sec:human-judgments}, to support our adoption of GPT-4 for evaluation purposes. Our results reveal a high concordance between GPT-4 and human judgments, with the alignment between human annotators and GPT-4 equaling or surpassing typical inter-annotator agreement.

\textbf{Methods.} Beyond DPO, we benchmark a range of established techniques for aligning language models with human preferences. As the most straightforward approaches, we investigate zero-shot prompting using \textbf{GPT-J} \citep{gpt-j} in the summarization experiments, and 2-shot prompting with \textbf{Pythia-2.8B} \citep{biderman2023pythia} for dialogue. We also assess the \textbf{SFT} model alongside the \textbf{Preferred-FT} variant, which is fine-tuned with supervised learning to produce the favored completion $y_w$—drawing training data either from the SFT model (for sentiment and summarization tasks) or a general language model (for dialogue). Another approach evaluated is \textbf{Unlikelihood}~\citep{welleck2019neural}, which optimizes the policy by increasing the probability of $y_w$ while simultaneously decreasing that of $y_l$, modulated by an optional coefficient $\alpha\in[0,1]$ for the unlikelihood penalty. Additionally, we include \textbf{PPO} \citep{schulman2017proximal}, leveraging a reward model trained on preference annotations, and \textbf{PPO-GT}, an oracle variant that directly learns from the ground-truth reward in the controlled sentiment case. For sentiment analysis, we test two PPO-GT implementations: a standard out-of-the-box version \cite{leandro_von_werra_2023_7790115} and a modified version incorporating reward normalization and advanced hyperparameter tuning for improved outcomes (these modifications are also applied to PPO with learned reward models). Finally, our evaluation includes the \textbf{Best of $N$} baseline, where $N$ responses are generated from either the SFT (or Preferred-FT in dialogue) model and the top response, as determined by a learned reward function, is selected. While this technique yields strong performance by separating reward model accuracy from PPO optimization, it is computationally prohibitive for even moderate $N$ since each test prompt requires sampling $N$ completions.

\subsection{How well can DPO optimize the RLHF objective?}

The KL-regularized reward maximization objective, standard in RLHF algorithms, serves to encourage reward maximization while simultaneously constraining the learned policy from straying too far from a reference policy. As a result, fair algorithmic comparisons should account for both the magnitude of attained reward and the extent of KL divergence; higher reward at the cost of substantially increased KL may not represent a meaningful improvement. Figure~\ref{fig:frontier-tldr-main} presents the tradeoff between reward and KL for several algorithms evaluated on the sentiment task. For each method, we perform a series of training experiments, each with a distinct hyperparameter governing policy conservatism (target KL $\in\{3,6,9,12\}$ for PPO, $\beta \in \{0.05,0.1,1,5\}$, $\alpha\in\{0.05,0.1,0.5,1\}$ for unlikelihood, and different random seeds for preferred-FT), resulting in a total of 22 runs. At intervals of every 100 training steps until reaching convergence, we assess each resulting policy using a fixed set of test prompts, calculating both the mean reward under the true reward function and the mean sequence-level KL\footnote{Specifically, this is computed as the sum of KL-divergences over all timesteps.} relative to the reference policy, $\text{KL}\left(\pi\mid \mid \pi_\text{ref}\right)$. The results show that DPO establishes a substantially more favorable reward-KL frontier, securing higher rewards with lower KL values compared to other approaches. This is noteworthy for several reasons. First, although DPO and PPO are designed to optimize the identical objective, DPO achieves markedly greater efficiency, with its reward/KL frontier strictly outperforming that of PPO. Second, DPO continues to outperform PPO even when PPO is supplied with ground-truth rewards (PPO-GT), yielding a superior reward-KL tradeoff in all evaluated regimes.

\subsection{Can DPO scale to real preference datasets?}
\label{sec:dpo-real-datasets}
We proceed to assess the fine-tuning capabilities of DPO on tasks involving summarization and single-turn dialogue. In summarization, it has been shown that automatic metrics like ROUGE do not always align closely with human judgment~\citep{stiennon2022learning}, and prior studies indicate that language models refined with PPO on human feedback tend to generate more desirable summaries. Our evaluation involves generating samples on the test partition of the TL;DR summarization dataset and measuring the mean win rate compared to reference summaries within the test set. For all approaches, outputs are sampled using temperatures from 0.0 up to 1.0, with corresponding win rates depicted in Figure~\ref{fig:frontier-tldr-main} (right). DPO, PPO, and Preferred-FT all fine-tune the same underlying GPT-J SFT model\footnote{\url{https://huggingface.co/CarperAI/openai_summarize_tldr_sft}}. Results reveal that DPO attains a win rate of about 61\% at a sampling temperature of 0.0, outperforming PPO, which reaches roughly 57\% at its optimal temperature setting. Furthermore, DPO achieves a greater peak win rate than the best-of-$N$ baseline. It is important to mention that DPO’s $\beta$ hyperparameter was not extensively optimized in these experiments, suggesting the reported outcomes may not fully reflect DPO’s optimal potential. Additionally, DPO demonstrates greater robustness to variations in sampling temperature relative to PPO, whose performance can decline to match that of the original GPT-J model at higher temperature settings. Preferred-FT offers little to no improvement over the SFT baseline. We further directly contrast DPO and PPO in human assessments in Section~\ref{sec:human-judgments}, where DPO-generated samples at temperature 0.25 were selected 58\% of the time over PPO-generated samples at temperature 0.

For the single-turn dialogue setting, our evaluation focuses on a subset of the Anthropic HH dataset's test split \citep{bai2022training}, specifically instances involving a single exchange between human and assistant. In the case of GPT-4, we determine win rates by referencing the preferred completions from the test set across the various methods. Since no established SFT baseline exists for this particular task, we initialize with the pre-trained Pythia-2.8B model. We then employ Preferred-FT to fine-tune a reference model on selected completions, ensuring that outputs remain representative of the model's distribution, before applying DPO training. Our comparisons also include the strongest result among 128 Preferred-FT completions (we observe that the Best of $N$ performance saturates at $N=128$; refer to Appendix Figure~\ref{fig:best-of-n}), alongside a 2-shot prompt version of the base Pythia-2.8B model. DPO is found to match or surpass these approaches at their respective optimal temperature settings. Furthermore, we assess an RLHF model trained with PPO on the Anthropic HH dataset\footnote{\url{https://huggingface.co/reciprocate/ppo_hh_pythia-6B}} sourced from a reputable repository\footnote{\url{https://github.com/CarperAI/trlx/tree/main/examples/hh}}, but we are unable to identify a prompting strategy or sampling temperature that consistently outperforms the unmodified Pythia-2.8B. Drawing on our TL;DR experiments and the fact that both PPO and Best of 128 target the same reward objective, we treat Best of 128 as an approximate benchmark for PPO-level results. In summary, DPO emerges as the only method that, while being computationally efficient, surpasses the preferred completions within the Anthropic HH dataset, achieving equal or better results relative to the much more resource-intensive Best of 128 method. Lastly, as depicted in Figure~\ref{fig:dialogue-main}, DPO rapidly converges to its optimal performance level.

\subsection{Generalization to a new input distribution}

To more thoroughly assess how PPO and DPO handle distributional shifts, we test the policies learned in our Reddit TL;DR summarization experiment on a distinct domain: news articles from the test split of the CNN/DailyMail dataset \citep{nallapati-etal-2016-abstractive}. The evaluation uses the optimal sampling temperatures determined from TL;DR (0 and 0.25). Results can be found in Table~\ref{tab:ood}. We determine the GPT-4 win rate by comparing system-generated summaries against ground-truth summaries from these datasets, applying the same GPT-4 (C) prompt as for Reddit TL;DR, with “forum post” swapped for “news article.” On this out-of-distribution data, DPO sustains a notable lead over PPO. These findings provide preliminary evidence that DPO is capable of generalizing to unseen distributions at least as effectively as PPO—even though DPO does not rely on the additional pool of unlabeled Reddit TL;DR prompts available to PPO.

\subsection{Validating GPT-4 judgments with human judgments}
\label{sec:human-judgments}
To assess the trustworthiness of GPT-4's evaluations, we carry out a human study leveraging results from the TL;DR summarization experiment and two distinct GPT-4 prompts. The \textbf{GPT-4 (S)} (simple) prompt requests a judgment on which summary better covers the post's essential content. In contrast, the \textbf{GPT-4 (C)} (concise) prompt asks not only about informativeness but also about conciseness, motivated by our observation that GPT-4, with the \textbf{GPT-4 (S)} prompt, shows a preference for longer, more repetitive summaries compared to human raters. Full prompt text is available in Appendix~\ref{app:prompts}. We select three system conditions for comparison: the top-performing method (DPO at temperature 0.25), the lowest-performing (PPO at temperature 1.0), and an intermediate system (SFT at temperature 0.25), ensuring a range of summary qualities; each is evaluated against PPO with greedy decoding (its best-performing temperature). Our results indicate that, using either prompt, GPT-4's assessments concur with human judgments at a rate comparable to the agreement among humans themselves (note that due to a small number of human raters, multiple ratings were only collected for DPO and PPO-1 comparisons). Notably, the \textbf{GPT-4 (C)} prompt tends to yield win rates that more closely match human preferences, so we adopt this prompt in our primary results in Section~\ref{sec:dpo-real-datasets}. Further details on the human evaluation process, including the rater interface and the roster of human volunteers, are provided in Appendix~\ref{app:human-study}.

\section{Discussion}
Preference-based learning offers a robust and scalable approach for developing capable and well-aligned language models. In this work, we have presented DPO, a straightforward training framework that leverages human preferences to train language models without employing reinforcement learning. DPO departs from traditional methods that recast preference learning as a standard reinforcement learning (RL) problem to take advantage of existing RL algorithms. Instead, DPO establishes a correspondence between language model policies and reward functions, allowing the model to be trained directly to align with human preferences using a simple cross-entropy loss, omitting the need for reinforcement learning and maintaining general applicability. Remarkably, with almost no hyperparameter tuning, DPO achieves performance on par with or superior to conventional RLHF approaches, such as those utilizing PPO. This substantially lowers the technical barriers to leveraging human preferences for language model training.

\textbf{Limitations \& Future Work.} Our findings prompt several avenues for further investigation. One pertinent question is how DPO policies handle distributional shifts in comparison to those trained using explicit reward functions. Preliminary evidence indicates that DPO policies may generalize comparably to those developed with PPO, yet this requires more systematic exploration. Additionally, it is worth exploring whether self-labeling with the DPO policy can efficiently leverage unlabeled prompts. Another important consideration is understanding how reward over-optimization presents itself in the direct preference optimization regime, particularly in light of the modest drop in performance observed in Figure~\ref{fig:dialogue-main}-right—whether this is an instance of such over-optimization remains to be determined. While our current evaluation focuses on models with up to 6B parameters, extending DPO to substantially larger, state-of-the-art models remains an enticing research direction. With regard to evaluation, we observed that the prompt can influence the win rates assessed by GPT-4; future studies could investigate optimal strategies for eliciting reliable automated evaluations. Ultimately, DPO’s potential applications are broad, extending well beyond human-preference-based language model training to include generative modeling in various other domains.